\documentclass[11pt]{article}
\usepackage[margin=1in]{geometry}
\usepackage{amsmath,amssymb,bm}
\usepackage{physics}
\usepackage{booktabs}

\title{RESTART: 1D Heat Equation with Hybrid CV-DV LCHS}
\date{\today}

\begin{document}
\maketitle

\section*{1. Objective and Source Alignment}
Goal: solve the 1D heat equation in the same operator target used by our finite-difference mapping,
\begin{equation}
\partial_t u(x,t)=\alpha\,\partial_x^2 u(x,t),\qquad u(0,t)=u(L,t)=0.
\end{equation}
For 4 interior points (2-qubit DV register),
\begin{equation}
\dot{\bm u}(t)=-\alpha T\bm u(t),\qquad
T=\begin{bmatrix}
2&-1&0&0\\
-1&2&-1&0\\
0&-1&2&-1\\
0&0&-1&2
\end{bmatrix}.
\end{equation}
The classical target is
\begin{equation}
\bm u_{\mathrm{PDE}}(t)=e^{-\alpha tT}\bm u(0).
\end{equation}

Reference baseline used in this restart:
\begin{itemize}
\item Working-paper draft: \texttt{res\_base/Hybrid\_CV\_DV\_LCHS (2).pdf} (Sec. 3.1, 3.5.2).
\item LCHS kernel theory: \texttt{res\_refs/Quantum algorithm for linear non-unitary dynamics with near-optimal dependence on all parameters.pdf}.
\item Constant-factor beta guidance: \texttt{res\_refs/Constant-Factor Improvements in Quantum Algorithms for Linear Differential Equations.pdf}.
\item Hybrid CV-DV gate conventions: \texttt{res\_refs/2407.10381v3.pdf} + installed \texttt{hybridlane}/\texttt{bosonic-qiskit} definitions.
\end{itemize}

\section*{2. LCHS Representation and Kernel Branch}
For homogeneous linear dynamics with $A=L+iH$ and $L\succeq 0$, LCHS writes
\begin{equation}
 e^{-At} = \int_{\mathbb R} g(k)\,e^{-it(kL+H)}\,dk,
\end{equation}
with
\begin{equation}
 g(k)=\frac{f(k)}{1-ik}.
\end{equation}
We use the improved-kernel family (parameter $\beta\in(0,1)$):
\begin{equation}
 f(k)=\frac{1}{C_\beta\,e^{(1+ik)^\beta}},\qquad C_\beta=2\pi e^{-2\beta},
\end{equation}
so
\begin{equation}
 g_\beta(k)=\frac{e^{-(1+ik)^\beta}}{C_\beta(1-ik)}.
\end{equation}
This matches the improved LCHS branch and gives faster truncation decay than the original Cauchy kernel.

\section*{3. CV State Engineering Used in Code}
Following the squeezed-state construction, coefficients are projected onto a squeezed Hermite/Fock basis:
\begin{equation}
C_n\propto \int dk\; H_n\!\left(\frac{k}{e^{r'}}\right)e^{-\gamma k^2}g_\beta(k),
\qquad \gamma=e^{-2r'}-e^{-2r},
\end{equation}
with required geometry
\begin{equation}
r'<r\quad\Longleftrightarrow\quad\gamma>0.
\end{equation}
The finite-$N$ truncation uses
\begin{equation}
\ket{\psi_{\mathrm{CV}}}\approx\sum_{n=0}^{N-1}C_n\ket{n},\qquad N=32.
\end{equation}
In our executable path we use incoherent Fock aggregation (mixture approximation):
\begin{equation}
\rho_{\mathrm{in}}\approx\sum_{n=0}^{N-1}|C_n|^2\ket n\bra n.
\end{equation}

\section*{4. Heat-Equation DV Generator and Circuit Terms}
For 2 qubits,
\begin{equation}
T=2(I\otimes I)-(I\otimes X)-\frac{1}{2}(X\otimes X+Y\otimes Y).
\end{equation}
Define coefficients
\begin{equation}
\lambda_I=2,\quad \lambda_{IX}=-1,\quad \lambda_{XX}=\lambda_{YY}=-\frac12.
\end{equation}
One first-order Trotter step with $\Delta t=t/n_{\text{steps}}$ implements
\begin{equation}
\exp[-i\Delta t\,\hat{k}\otimes (\lambda_I I + \lambda_{IX}IX + \lambda_{XX}XX + \lambda_{YY}YY)]
\end{equation}
via basis changes + controlled displacements.

\section*{5. Hybridlane Gate Mapping and Branch Calibration}
Installed gate definition:
\begin{equation}
D(\alpha)=\exp(\alpha a^\dagger-\alpha^*a),\qquad \alpha=a_0e^{i\phi}.
\end{equation}
The displacement phase therefore selects a quadrature branch for the effective $\hat{k}$ coupling.
A historical failure mode was assuming this branch incorrectly from notation only.

In RESTART, we therefore treat phase as an explicit model parameter and provide a built-in calibration test:
\begin{equation}
\phi\in\left\{0,-\frac\pi2,+\frac\pi2\right\},
\end{equation}
ranked by primary metric $\varepsilon_{\mathrm{PDE}}$ (below).

\section*{6. Metrics and Error Budget}
Primary metric (project objective):
\begin{equation}
\varepsilon_{\mathrm{PDE}}=
\frac{\norm{\bm u_{\mathrm{PDE}}-\sqrt{p_{\mathrm{post}}}\,\psi_{\mathrm{principal}}}}{\norm{\bm u_{\mathrm{PDE}}}}.
\end{equation}
Here $\psi_{\mathrm{principal}}$ is the principal eigenvector of the reconstructed post-selected DV density matrix.

Mixed-state directional diagnostic:
\begin{equation}
F=\mel{\hat u_{\mathrm{PDE}}}{\rho_{\mathrm{post}}}{\hat u_{\mathrm{PDE}}},
\qquad \hat u_{\mathrm{PDE}}=\frac{\bm u_{\mathrm{PDE}}}{\norm{\bm u_{\mathrm{PDE}}}}.
\end{equation}

Practical decomposition of total mismatch:
\begin{equation}
\varepsilon_{\mathrm{total}}\lesssim
\varepsilon_{\mathrm{alg}}+\varepsilon_{\mathrm{trotter}}+\varepsilon_{\mathrm{cutoff}}+\varepsilon_{\mathrm{mix}}+\varepsilon_{\mathrm{num}}.
\end{equation}
In this setup, $\varepsilon_{\mathrm{alg}}=0$ by construction (same target operator $T$ in theory and classical reference).
A first-order Lie--Trotter bound scales as
\begin{equation}
\varepsilon_{\mathrm{trotter}}=\mathcal O\!\left(\frac{(\alpha t)^2}{n_{\text{steps}}}\right),
\end{equation}
with explicit commutator constant computed in code.

\section*{7. Theorem-Constrained Sensitivity (No Random Search)}
Search variables:
\begin{equation}
\theta=(r,r',\beta),
\end{equation}
with deterministic grids and local deterministic refinement.

Hard/theory constraints used:
\begin{align}
&r'<r-\Delta_{\min},\\
&\gamma=e^{-2r'}-e^{-2r}\ge \gamma_{\min},\\
&\beta\in(0,1),\\
&M_{\mathrm{tail}}(N)\le M_{\max},\\
&n_{\mathrm{eff}}=\frac{1}{\sum_n|C_n|^4}\le n_{\mathrm{eff,max}}.
\end{align}

Ranking objective (PDE-first):
\begin{equation}
J(\theta)=\varepsilon_{\mathrm{PDE}}(\theta)
+\lambda_p\frac{[p_{\min}-p_{\mathrm{post}}(\theta)]_+}{p_{\min}}
-\eta\,p_{\mathrm{post}}(\theta),
\end{equation}
with feasible points prioritized over infeasible ones.

\section*{8. Implementation Files (RESTART Folder)}
\begin{itemize}
\item \texttt{RESTART/heat1d\_lchs\_cv\_dv.py}: single-run circuit + metrics + phase calibration.
\item \texttt{RESTART/heat1d\_lchs\_sensitivity.py}: theorem-constrained deterministic sensitivity/optimization.
\end{itemize}

\section*{9. Reproducible Commands}
Single run:
\begin{verbatim}
python RESTART/heat1d_lchs_cv_dv.py --calibrate-phase
\end{verbatim}

Sensitivity run:
\begin{verbatim}
python RESTART/heat1d_lchs_sensitivity.py \
  --output-dir RESTART/results \
  --n-r-target 5 --n-r-prime 5 --n-beta 7
\end{verbatim}

\end{document}
